% Options for packages loaded elsewhere
\PassOptionsToPackage{unicode}{hyperref}
\PassOptionsToPackage{hyphens}{url}
\documentclass[
  ignorenonframetext,
]{beamer}
\newif\ifbibliography
\usepackage{pgfpages}
\setbeamertemplate{caption}[numbered]
\setbeamertemplate{caption label separator}{: }
\setbeamercolor{caption name}{fg=normal text.fg}
\beamertemplatenavigationsymbolsempty
% remove section numbering
\setbeamertemplate{part page}{
  \centering
  \begin{beamercolorbox}[sep=16pt,center]{part title}
    \usebeamerfont{part title}\insertpart\par
  \end{beamercolorbox}
}
\setbeamertemplate{section page}{
  \centering
  \begin{beamercolorbox}[sep=12pt,center]{section title}
    \usebeamerfont{section title}\insertsection\par
  \end{beamercolorbox}
}
\setbeamertemplate{subsection page}{
  \centering
  \begin{beamercolorbox}[sep=8pt,center]{subsection title}
    \usebeamerfont{subsection title}\insertsubsection\par
  \end{beamercolorbox}
}
% Prevent slide breaks in the middle of a paragraph
\widowpenalties 1 10000
\raggedbottom
\AtBeginPart{
  \frame{\partpage}
}
\AtBeginSection{
  \ifbibliography
  \else
    \frame{\sectionpage}
  \fi
}
\AtBeginSubsection{
  \frame{\subsectionpage}
}
\usepackage{iftex}
\ifPDFTeX
  \usepackage[T1]{fontenc}
  \usepackage[utf8]{inputenc}
  \usepackage{textcomp} % provide euro and other symbols
\else % if luatex or xetex
  \usepackage{unicode-math} % this also loads fontspec
  \defaultfontfeatures{Scale=MatchLowercase}
  \defaultfontfeatures[\rmfamily]{Ligatures=TeX,Scale=1}
\fi
\usepackage{lmodern}
\ifPDFTeX\else
  % xetex/luatex font selection
\fi
% Use upquote if available, for straight quotes in verbatim environments
\IfFileExists{upquote.sty}{\usepackage{upquote}}{}
\IfFileExists{microtype.sty}{% use microtype if available
  \usepackage[]{microtype}
  \UseMicrotypeSet[protrusion]{basicmath} % disable protrusion for tt fonts
}{}
\makeatletter
\@ifundefined{KOMAClassName}{% if non-KOMA class
  \IfFileExists{parskip.sty}{%
    \usepackage{parskip}
  }{% else
    \setlength{\parindent}{0pt}
    \setlength{\parskip}{6pt plus 2pt minus 1pt}}
}{% if KOMA class
  \KOMAoptions{parskip=half}}
\makeatother
\usepackage{longtable,booktabs,array}
\newcounter{none} % for unnumbered tables
\usepackage{calc} % for calculating minipage widths
\usepackage{caption}
% Make caption package work with longtable
\makeatletter
\def\fnum@table{\tablename~\thetable}
\makeatother
\usepackage{graphicx}
\makeatletter
\newsavebox\pandoc@box
\newcommand*\pandocbounded[1]{% scales image to fit in text height/width
  \sbox\pandoc@box{#1}%
  \Gscale@div\@tempa{\textheight}{\dimexpr\ht\pandoc@box+\dp\pandoc@box\relax}%
  \Gscale@div\@tempb{\linewidth}{\wd\pandoc@box}%
  \ifdim\@tempb\p@<\@tempa\p@\let\@tempa\@tempb\fi% select the smaller of both
  \ifdim\@tempa\p@<\p@\scalebox{\@tempa}{\usebox\pandoc@box}%
  \else\usebox{\pandoc@box}%
  \fi%
}
% Set default figure placement to htbp
\def\fps@figure{htbp}
\makeatother
\setlength{\emergencystretch}{3em} % prevent overfull lines
\providecommand{\tightlist}{%
  \setlength{\itemsep}{0pt}\setlength{\parskip}{0pt}}
\usepackage{bookmark}
\IfFileExists{xurl.sty}{\usepackage{xurl}}{} % add URL line breaks if available
\urlstyle{same}
\hypersetup{
  hidelinks,
  pdfcreator={LaTeX via pandoc}}

\author{\texorpdfstring{}{}}
\date{}

\begin{document}

\begin{frame}
\newpage
\end{frame}

\begin{frame}[fragile]{Attack Corpus}
\protect\phantomsection\label{sec:attack-corpus}
This supplementary material provides corpus overview
(\cref{sec:corpus-overview}), detailed statistics
(\cref{sec:corpus-stats}), example attacks by category
(\cref{sec:attack-examples}), generation methodology
(\cref{sec:generation-methodology}), effectiveness analysis
(\cref{sec:effectiveness-analysis}), and ethical considerations
(\cref{sec:ethical-considerations}).

\begin{block}{Corpus Overview}
\protect\phantomsection\label{sec:corpus-overview}
The attack corpus used for experimental validation comprises 950 unique
attack instances across four primary categories. This supplementary
material provides detailed statistics, sanitized examples, generation
methodology, and ethical considerations.

The attack corpus is generated programmatically via deterministic random
seeds (seed=42 for all reported results). The attack taxonomy is defined
as a Python enum (\texttt{AttackCategory} in
\texttt{src/utils/types.py}) with 4 top-level categories and 12
subcategories. No static data file is shipped; the corpus is regenerated
on each evaluation run via \texttt{AttackCorpus.generate()} to ensure
reproducibility. Researchers can serialize the corpus to JSON via
\texttt{AttackCorpus.save()} for inspection.
\end{block}

\begin{block}{Full Attack Corpus Statistics}
\protect\phantomsection\label{sec:corpus-stats}
\begin{figure}
\centering
\includegraphics[width=0.95\linewidth,height=\textheight,keepaspectratio,alt={Cognitive Attack Taxonomy. Hierarchical visualization of the 950-attack corpus organized by primary category (Prompt Injection, Trust Exploitation, Belief Manipulation, Coordination Attacks) and subcategory. Node size indicates attack count; color intensity indicates baseline success rate against undefended systems. The taxonomy classifies attacks by five adversary capability classes (Part 1, Definition 4): \textbackslash Omega\_1 (passive eavesdropping), \textbackslash Omega\_2 (message injection), \textbackslash Omega\_3 (identity spoofing), \textbackslash Omega\_4 (belief manipulation), and \textbackslash Omega\_5 (coordinated multi-agent attacks). Prompt injection dominates in volume (500 attacks, 53\% of corpus) while coordination attacks show highest baseline success rate (82\%) due to their ability to exploit the absence of inter-agent verification. Generated deterministically with seed 42 via .}]{../figures/comprehensive_taxonomy.pdf}
\caption{Cognitive Attack Taxonomy. Hierarchical visualization of the
950-attack corpus organized by primary category (Prompt Injection, Trust
Exploitation, Belief Manipulation, Coordination Attacks) and
subcategory. Node size indicates attack count; color intensity indicates
baseline success rate against undefended systems. The taxonomy
classifies attacks by five adversary capability classes (Part 1,
Definition 4): \(\Omega_1\) (passive eavesdropping), \(\Omega_2\)
(message injection), \(\Omega_3\) (identity spoofing), \(\Omega_4\)
(belief manipulation), and \(\Omega_5\) (coordinated multi-agent
attacks). Prompt injection dominates in volume (500 attacks, 53\% of
corpus) while coordination attacks show highest baseline success rate
(82\%) due to their ability to exploit the absence of inter-agent
verification. Generated deterministically with seed 42 via
\texttt{AttackCorpus.generate()}.}\label{fig:comprehensive-taxonomy}
\end{figure}

The cognitive attack taxonomy (\cref{fig:comprehensive-taxonomy})
organizes our 950-attack corpus into a hierarchical structure that
reflects both the attack mechanisms and their relative prevalence in the
wild.

\begin{block}{Category Breakdown}
\protect\phantomsection\label{sec:category-breakdown}
\textbf{Table: Attack corpus composition by category.}
\{\#tab:corpus-categories\}

{\def\LTcaptype{none} % do not increment counter
\begin{longtable}[]{@{}lllll@{}}
\toprule\noalign{}
Category & Total & Train & Test & Validation \\
\midrule\noalign{}
\endhead
Prompt Injection & 500 & 350 & 100 & 50 \\
Trust Exploitation & 200 & 140 & 40 & 20 \\
Belief Manipulation & 150 & 105 & 30 & 15 \\
Coordination Attacks & 100 & 70 & 20 & 10 \\
\textbf{Total} & \textbf{950} & \textbf{665} & \textbf{190} &
\textbf{95} \\
\bottomrule\noalign{}
\end{longtable}
}
\end{block}

\begin{block}{Prompt Injection Subcategories}
\protect\phantomsection\label{sec:injection-subcats}
\textbf{Table: Prompt injection subcategory statistics.}
\{\#tab:injection-subcats\}

{\def\LTcaptype{none} % do not increment counter
\begin{longtable}[]{@{}llll@{}}
\toprule\noalign{}
Subcategory & Count & Undefended Success Rate\(^*\) & CIF Success \\
\midrule\noalign{}
\endhead
Direct injection & 200 & 78\% & 3\% \\
Indirect injection & 200 & 65\% & 5\% \\
Nested injection & 100 & 82\% & 7\% \\
\bottomrule\noalign{}
\end{longtable}
}

\(^*\)\textit{Undefended success rates represent attack effectiveness measured without any CIF defense mechanisms active.}

\textbf{Direct Injection}: Attacks embedded directly in user input
attempting to override system instructions.

\textbf{Indirect Injection}: Attacks injected through external data
sources (web content, API responses, documents).

\textbf{Nested Injection}: Multi-layer attacks where outer content masks
inner malicious payloads.
\end{block}

\begin{block}{Trust Exploitation Subcategories}
\protect\phantomsection\label{sec:trust-subcats}
\textbf{Table: Trust exploitation subcategory statistics.}
\{\#tab:trust-subcats\}

{\def\LTcaptype{none} % do not increment counter
\begin{longtable}[]{@{}lll@{}}
\toprule\noalign{}
Subcategory & Count & Description \\
\midrule\noalign{}
\endhead
Identity impersonation & 80 & Claiming to be trusted entity \\
Trust inflation & 70 & Artificially boosting trust scores \\
Delegation abuse & 50 & Exploiting delegation chains \\
\bottomrule\noalign{}
\end{longtable}
}
\end{block}

\begin{block}{Belief Manipulation Subcategories}
\protect\phantomsection\label{sec:belief-subcats}
\textbf{Table: Belief manipulation subcategory statistics.}
\{\#tab:belief-subcats\}

{\def\LTcaptype{none} % do not increment counter
\begin{longtable}[]{@{}lll@{}}
\toprule\noalign{}
Subcategory & Count & Description \\
\midrule\noalign{}
\endhead
Direct belief injection & 60 & Asserting false facts \\
Evidence fabrication & 50 & Creating fake supporting evidence \\
Progressive drift & 40 & Gradual belief modification \\
\bottomrule\noalign{}
\end{longtable}
}
\end{block}

\begin{block}{Coordination Attack Subcategories}
\protect\phantomsection\label{sec:coord-subcats}
\textbf{Table: Coordination attack subcategory statistics.}
\{\#tab:coord-subcats\}

{\def\LTcaptype{none} % do not increment counter
\begin{longtable}[]{@{}lll@{}}
\toprule\noalign{}
Subcategory & Count & Description \\
\midrule\noalign{}
\endhead
Sybil attacks & 40 & Fake agent injection \\
Consensus poisoning & 35 & Corrupting multi-agent agreement \\
Timing attacks & 25 & Exploiting synchronization \\
\bottomrule\noalign{}
\end{longtable}
}
\end{block}

\begin{block}{Detailed Statistics by Source}
\protect\phantomsection\label{sec:source-stats}
\textbf{Table: Attack source distribution.} \{\#tab:attack-sources\}

{\def\LTcaptype{none} % do not increment counter
\begin{longtable}[]{@{}lll@{}}
\toprule\noalign{}
Source & Count & Percentage \\
\midrule\noalign{}
\endhead
Published datasets & 320 & 33.7\% \\
Red team exercises & 280 & 29.5\% \\
Synthetic generation & 200 & 21.1\% \\
Custom adversarial & 150 & 15.8\% \\
\bottomrule\noalign{}
\end{longtable}
}

\textbf{Table: Published dataset sources.} \{\#tab:dataset-sources\}

{\def\LTcaptype{none} % do not increment counter
\begin{longtable}[]{@{}lll@{}}
\toprule\noalign{}
Dataset & Attacks Used & Citation \\
\midrule\noalign{}
\endhead
JailbreakBench & 150 & \cite{chao2024jailbreakbench} \\
PromptInject & 80 & \cite{liu2023prompt} \\
TensorTrust & 50 & \cite{toyer2024tensortrust} \\
Custom academic & 40 & Various \\
\bottomrule\noalign{}
\end{longtable}
}
\end{block}

\begin{block}{Complexity Distribution}
\protect\phantomsection\label{sec:complexity-dist}
\textbf{Table: Attack complexity distribution.}
\{\#tab:complexity-dist\}

{\def\LTcaptype{none} % do not increment counter
\begin{longtable}[]{@{}llll@{}}
\toprule\noalign{}
Complexity Level & Count & Average Tokens & Detection Difficulty \\
\midrule\noalign{}
\endhead
Low & 250 & 45 & Easy \\
Medium & 400 & 120 & Moderate \\
High & 200 & 280 & Hard \\
Adversarial & 100 & 450 & Expert \\
\bottomrule\noalign{}
\end{longtable}
}
\end{block}

\begin{block}{Target Distribution}
\protect\phantomsection\label{sec:target-dist}
\textbf{Table: Attack target distribution.} \{\#tab:target-dist\}

{\def\LTcaptype{none} % do not increment counter
\begin{longtable}[]{@{}lll@{}}
\toprule\noalign{}
Target & Count & Category \\
\midrule\noalign{}
\endhead
Belief state & 280 & Epistemic \\
Action execution & 250 & Behavioral \\
Trust relationships & 220 & Social \\
Temporal state & 100 & Persistence \\
Goal alignment & 100 & Behavioral \\
\bottomrule\noalign{}
\end{longtable}
}

\begin{figure}
\centering
\includegraphics[width=0.9\linewidth,height=\textheight,keepaspectratio,alt={Attack Surface Map. Visualization of cognitive attack entry points in multiagent systems showing the five primary attack surfaces aligned to CIF defense layers: User Input (direct injection, countered by Cognitive Firewall), Tool Outputs (indirect injection, countered by Belief Sandbox + Firewall), Agent Communication (trust exploitation via delegation chains, countered by Trust Calculus with \textbackslash delta\^{}d decay), Persistent Memory (belief poisoning and progressive drift, countered by Tripwires + Drift Detection), and External Triggers (timing and coordination attacks, countered by Byzantine Consensus). Line thickness indicates attack frequency in our 950-attack corpus; node color indicates CIF detection efficacy at each surface (green \$\textgreater\$90\%, yellow 80-90\%, red \$\textless\$80\%). The diagram highlights that Agent Communication is both the highest-frequency and highest-risk surface in multiagent deployments---a gap that single-agent defenses do not address.}]{../figures/attack_surface.pdf}
\caption{Attack Surface Map. Visualization of cognitive attack entry
points in multiagent systems showing the five primary attack surfaces
aligned to CIF defense layers: User Input (direct injection, countered
by Cognitive Firewall), Tool Outputs (indirect injection, countered by
Belief Sandbox + Firewall), Agent Communication (trust exploitation via
delegation chains, countered by Trust Calculus with \(\delta^d\) decay),
Persistent Memory (belief poisoning and progressive drift, countered by
Tripwires + Drift Detection), and External Triggers (timing and
coordination attacks, countered by Byzantine Consensus). Line thickness
indicates attack frequency in our 950-attack corpus; node color
indicates CIF detection efficacy at each surface (green
\$\textgreater\$90\%, yellow 80-90\%, red \$\textless\$80\%). The
diagram highlights that Agent Communication is both the
highest-frequency and highest-risk surface in multiagent deployments---a
gap that single-agent defenses do not
address.}\label{fig:attack-surface}
\end{figure}

The attack surface map (\cref{fig:attack-surface}) visualizes how these
attack categories map to distinct entry points in multiagent system
architectures.
\end{block}
\end{block}

\begin{block}{Detailed Attack Content}
\protect\phantomsection\label{detailed-attack-content}
The following subsections provide detailed attack examples, methodology,
and ethical considerations:

\begin{itemize}
\item
  \textbf{Attack Examples} (\cref{sec:attack-examples}): Sanitized
  examples for all four categories --- prompt injection (direct,
  indirect, nested), trust exploitation (impersonation, inflation,
  delegation abuse), belief manipulation (injection, fabrication,
  drift), and coordination attacks (Sybil, consensus poisoning, timing).
  See \texttt{03b\_attack\_examples.md}.
\item
  \textbf{Methodology and Ethics} (\cref{sec:generation-methodology}):
  Attack generation methodology, effectiveness analysis, ethical
  considerations (responsible disclosure, dual-use, human subjects), and
  data availability. See \texttt{03c\_attack\_ethics.md}.
\end{itemize}
\end{block}
\end{frame}

\end{document}
